\item Comience con la ecuación 6.11 y demuestre que la presión total $P$ en un recipiente lleno con una mezcla de varios gases ideales es $P = P_1 + P_2 + P_3 + \dots$, donde $P_1, P_2, \dots$ son las presiones que cada gas ejercería si llenara solo el recipiente. (Estas presiones individuales se llaman \textit{presiones parciales} de los gases respectivos.) Este resultado se conoce como \textit{ley de Dalton de presiones parciales}.
